\section{Limitations and Evaluation Gaps}
\label{sec:limitations}

In accordance with rigorous academic integrity and the evaluation protocol established in our audit, this section explicitly delineates the current empirical limitations of our evaluation and presents the planned benchmark protocol.

\subsection{Absence of Large-Scale Benchmark Ground-Truth Evaluation}
While SecretGuard's unit test suite (112 test cases) and micro-benchmarks confirm that the mathematical scoring model successfully discriminates test fixtures from production credentials on isolated samples, the system has not yet been systematically evaluated against large-scale public ground-truth corpora such as \texttt{SecretBench} (Basak et al.~\cite{basak2023comparative}). Consequently, while marketing claims of ``near-zero false alarms'' exist in product documentation, we explicitly refrain from asserting empirical F1-score superiority over mature static scanners such as Gitleaks or TruffleHog until such comparative trials are conducted.

\subsection{Proposed Future Empirical Benchmark Protocol}
To address this evaluation gap, we outline a concrete, reproducible evaluation protocol to be executed in subsequent studies:
\begin{enumerate}
    \item \textbf{Ground-Truth Dataset Ingestion:} Evaluate SecretGuard against the complete \texttt{SecretBench} and \texttt{FPSecretBench} repositories~\cite{basak2023comparative}, comprising thousands of validated true-positive and false-positive instances across diverse languages.
    \item \textbf{Comparative Baselines:} Benchmark SecretGuard in two configurations---Deterministic Offline Mode ($S_{amb} < 0.20$ rules only) and Full Adaptive Mode (with Profile A and Profile B LLM verification)---directly against Gitleaks v8.18+, TruffleHog v3.60+, and Yelp \texttt{detect-secrets} v1.4+.
    \item \textbf{Statistical Metrics:} Compute Precision, Recall, and F1-score alongside the ratio of network calls $\rho_{net} = \frac{N_{\text{LLM}}}{N_{\text{Total}}}$ to quantify the exact reduction in false alarms and cloud invocation overhead.
\end{enumerate}

\subsection{Scope of Automated In-Place Refactoring}
The current automated quarantine engine operates on regular expression AST heuristics for Python and JavaScript/TypeScript files. In scenarios involving complex multi-line string interpolations, nested dictionary configurations, or languages with mandatory error handling (e.g., Go's \texttt{os.Getenv} or Rust's \texttt{std::env::var().unwrap()}), automated in-place replacement requires AST synthesis to avoid introducing syntax or compilation errors. Expanding the remediation engine to robust compiler-backed AST transformations is planned for future work.
